\documentclass[a4paper]{article}
\usepackage[T1]{fontenc}
\usepackage{lmodern}
\usepackage[margin=22mm]{geometry}
\usepackage[code,calculations,diagrams,requirements]{pxCORE}
\title{P3CTeX: reusable student tools}
\author{P3CTeX}
\date{Development guide, October 2026}
\begin{document}
\maketitle
\section{Choose the tools you need}
\begin{verbatim}
\documentclass[code,images,tables,diagrams,
               calculations,requirements]{P3CTeX}
\end{verbatim}
Existing module options remain supported. The diagram entry point
\texttt{pxDIAG} uses the established \texttt{pxCMP} engine and presets.
The new calculation package is \texttt{pxCALC}. Package names do not
dictate the keywords used to load tools.

\section{Write a formula once}
\begin{verbatim}
\pxDefVar{m}{1200}[m]
\pxDefVar{v}{25}[v]
\pxDefFormula{Ek}{\frac{1}{2}[m][v]^2}
  [symbol=E_k,unit=J]
\pxSolve{Ek}
\pxSolve[values={m=800,v=40}]{Ek}
\end{verbatim}
\pxDefVar{m}{1200}[m]
\pxDefVar{v}{25}[v]
\pxDefFormula{Ek}{\frac{1}{2}[m][v]^2}[symbol=E_k,unit=J]
\[\pxSolve{Ek}\]
\[\pxSolve[values={m=800,v=40}]{Ek}\]
The override is local. The next call still uses the original values.
\texttt{steps=1} prints the result; \texttt{steps=2} also prints the formula;
\texttt{steps=3} also prints substituted values. Settings can be supplied
for each call; formula definitions store only their symbol and unit.

\verb|\pxCalcStore{Ek}{\energy}| saves the unrounded number for reuse.
\verb|\pxSetVar{v}{30}| changes the value in the current TeX group.
\verb|\pxCalc{sin(pi/2)}| evaluates a numerical expression directly.
Units format results; they do not convert values. Divide joules by 1000
before labelling the result as kilojoules.

\section{Verbatim terminal and source ranges}
\begin{pxTerminal}{Shell session}
$ printf '%s\n' "$HOME" '100%'
/home/student
100%
\end{pxTerminal}
\begin{verbatim}
\begin{pxTerminal}{Shell session}
$ printf '%s\n' "$HOME" '100%'
/home/student
100%
\end{pxTerminal}
\pxCodeInline[pxBash]|echo "$HOME" # literal code|
\pxCodeChunk[pxJava]{Main.java}{12}{20}{Validation}
\end{verbatim}
Terminal options use listings keys, like \texttt{pxCode}. Palettes and
language definitions remain customisable through the existing listing API.
Inline verbatim must be outside other commands' arguments.

\section{Requirements and analysis}
\begin{pxUserStory}{US-102: Reuse calculations}
\pxAsA{student}
\pxIWant{to change inputs without copying the formula}
\pxSoThat{the displayed working stays consistent}
\pxAcceptanceCriteria{Show substitutions and preserve the defaults.}
\end{pxUserStory}
\begin{pxScenario}{Change one input}
\pxGiven{a speed of 25 metres per second}
\pxWhen{I solve with a speed of 30}
\pxThen{the result is 540\,000 joules}
\end{pxScenario}
\pxSWOT{Reusable formulas}{Explicit supported syntax}
  {Student laboratory reports}{Incorrect input data}

Cards accept normal \texttt{tcolorbox} options; their labels are editable
with \verb|\pxREQsetup{labels={as-a=Com a,i-want=Vull}}|.
The SWOT grid reuses \texttt{pxTAB}; its optional argument accepts table
options, including an explicit column specification.

\section{Reuse the existing environment}
Safe images, placeholders and image lists already belong to \texttt{pxSRC}.
Transition tables already include partition labels in \texttt{pxCMP}.
Use \texttt{pxMEM} for report annexes and the class for cover metadata and
authorship declarations. These tools do not introduce replacement systems.

The companion \texttt{student-tools.md} gives the complete syntax limits,
customisation keys and migration details.
\end{document}
